\scanpage{94}
$Q'$ : clôture \unsure{tot$^{\rm t}$|tot.} réelle de $Q$ ....

\margin{$\alpha=2$}
% circled
La condition a équivaut à la suivante, si $q\in Q^{* \unsure{-|{}}}$ est $>0$ :
\diagram{\begin{tikzcd}[ampersand replacement=\&,row sep=2pt]
\parbox[t]{.87\linewidth}{$\bigg\{$ \deleted{\unsure{on a $M$ \illegible{1}}}
\inserted{Pour tout $\lambda\in M$, $\exists$ \unsure{$r|m$}$>0$ tel que}
$\deleted{\unsure{\lambda}}\lambda^{r}\in Q'(i)^*$, et pour tout $\lambda\in M$, il existe un
entier $i$ tel que $|\lambda|=q^{i/2}$ i.e. $\lambda\bar\lambda=q^i$}\\
\Updownarrow\\
\parbox[t]{.87\linewidth}{$\bigg\{$ Pour tout $\lambda\in M$, \deleted{il existe} on a $|\lambda|=q^{j^*(\lambda)/2}$ i.e. $\lambda\bar\lambda=q^{j^*(\lambda)}$}
\end{tikzcd}}
\margin{N.B. Le $i$ en question alors est uniquement déterminé et n'est autre que $j^*(\lambda)$.}

\uline{Cor} Soit $\lambda\in\bar Q$, tel que \deleted{il existe $q\in Q^{\unsure{\times|*}+}$, avec}
\[
\lambda_\alpha\bar\lambda_\alpha=\lambda\bar\lambda\boxed{\in Q}
\quad\text{pour tout conjugué $\lambda_\alpha$ de $\lambda$. Alors $\exists\unsure{r|n}>0$,}
\]
\inserted{$\unsure{\Vert|=}\ q$}\quad \inserted{$|\lambda_\alpha|^2=|\lambda|^2$}
tel que $\lambda^{r}\in Q'(i)$, [donc la condition ne dépend pas du
plongt de $Q(\lambda)$ dans $\bar Q\subset\unsure{C|\mathbb C}$], et \unsure{.?.}
il existe $M,\varepsilon,j\ldots$ \inserted{relatifs à $f$} tel que $\lambda\in M$ (et \unsure{réciprt|vice versa}).

\margin{\uline{N.B.} Si un tel $\lambda$ est réel (donc tot réel !) on a $\lambda^2=q^i$ donc $\lambda=(q^{1/2})^i$, \uline{cas trivial}}
% slanted
\uline{Proposition} Soit $\lambda\in\bar Q^*$, et $q\in Q^{*+}$. Conditions équivalentes

(i) $\exists$ entier $i$ tel que pour $\forall$ conjugué $\lambda_\alpha$ de $\lambda$, on ait $|\lambda_\alpha|=q^{i/2}$
i.e. $\lambda_\alpha\bar\lambda_\alpha=q^i$

(ii bis) $\exists$ entier $n>0$ tel que $\lambda^n\in Q'(i)$, et \deleted{\unsure{$\exists i$|$\exists$}} entier $i$ tel que
$|\lambda|=q^{i/2}$ i.e. $\lambda\bar\lambda=q^i$ i.e. $\lambda^n\overline{\lambda^n}=q^{in}$

(ii) $\lambda\in Q'(i)$, il existe entier $i$ tel que $|\lambda|=q^{i/2}$, \unsure{id.|i.e.} $\lambda\bar\lambda=q^i$

(iii) $\exists$ un groupe de t.m. $G$ sur $Q$, et un élément $f$ de $G(Q)$,
$\deleted{\text{tel que l'on}}\ G_m\mathrel{\mathop{-}\limits^{j}}G\mathrel{\mathop{-}\limits^\varepsilon}G_m$, avec $\varepsilon j(\lambda)=\lambda^2$,
$\varepsilon(f)=q$, \inserted{avec et $\operatorname{Ker}\varepsilon$ compact}, tels que $\lambda$ soit une valeur propre de $f$
pour un $G$-Module de dim finie \inserted{sur $Q$} \unsure{convenable}, i.e.
$\lambda$ appartient à l'image de l'hom $M=D(G_{\bar Q})\longrightarrow\bar Q^*$
défini par $f$.

\margin{\deleted{NB. Les conditions \uline{n'impliquent pas} que $\lambda\in Q'(i)^*$, comme on voit p. ex. en prenant $\lambda=p^{1/4}$ ; Alors $Q(\lambda)$ n'est ni tot. imaginaire ($Q(p^{1/4})$ est réel) ni tot réel (car il est conjugué au $Q(\sqrt{\unsure{-p^{1/2}}})$ qui est \unsure{imaginaire}, donc n'est pas $\subset Q'(i)$}}
% slanted; crossed with long near-vertical strokes
(iv) Il existe une algèbre de Riemann commutative \deleted{\unsure{ab}}$A$ sur $Q$ et $f\in A^{\unsure{*|1|{}}}$
\inserted{et $i\in\mathbb N$, tels que $ff'=q^{i}$,}
et un hom $A\deleted{\unsure{C}}\longrightarrow\bar Q$, appliquant $f$ dans $\lambda$.
% archive sheet number: 93

